% Scientific Data manuscript draft: References section
% Compile from this directory:
%   pdflatex Scientific_Data_References_Draft_2026-07-22.tex
%   pdflatex Scientific_Data_References_Draft_2026-07-22.tex
\documentclass[11pt,letterpaper]{article}

\usepackage[margin=1in]{geometry}
\usepackage{booktabs}
\usepackage{longtable}
\usepackage{hyperref}
\usepackage{fancyhdr}
\usepackage{enumitem}
\usepackage{xcolor}
\usepackage[T1]{fontenc}
\usepackage[utf8]{inputenc}
\usepackage{lmodern}
\setlength{\headheight}{14pt}

\hypersetup{
  colorlinks=true,
  linkcolor=black,
  urlcolor=blue,
  citecolor=black,
  pdftitle={Scientific Data draft: References},
  pdfauthor={Alex Rees - Shmuel Lab}
}

\pagestyle{fancy}
\fancyhf{}
\fancyhead[L]{Scientific Data draft -- References}
\fancyhead[R]{Internal -- verify before submission}
\fancyfoot[C]{\thepage}
\renewcommand{\headrulewidth}{0.4pt}

\title{Scientific Data Data Descriptor\\[0.35em]
\large Draft: References\\[0.2em]
\large Software inventory and key literature}
\author{Prepared from project software versions and verified bibliographic sources\\
\small Alex Rees -- Shmuel Lab / \texttt{neuro\_pipeline}}
\date{22 July 2026}

\begin{document}
\maketitle

\begin{abstract}
\noindent
This draft provides a \textbf{References} section suitable for a Scientific Data
Data Descriptor, focused on the software and literature actually used in the
present workflow (BIDS conversion, de-identification, MRIQC, Pizarro screening,
defacing, and validation).
Peer-reviewed citations below were checked against publisher/DOI records or
the software authors' own citation instructions.
Where only a software version/URL is known, the item is listed under
\textit{Software and standards} rather than given an invented journal citation.
\end{abstract}

\tableofcontents
\newpage

\section*{Author notes (not for submission)}

\begin{itemize}[leftmargin=*]
  \item Pizarro citation text follows the authors' GitHub README for
        \texttt{AS-Lab/Pizarro-et-al-2023-DL-detects-MRI-artifacts}.
  \item MRIQC version recorded in derivatives metadata is
        \texttt{24.1.0.dev0+gd5b13cb5.d20240826}; array scripts also reference
        container tag \texttt{nipreps/mriqc:24.0.2}. Cite the MRIQC paper and
        report the exact runtime version in Methods.
  \item bids-validator runtime used for Scientific Data validation reports was
        \textbf{1.15.0}; pipeline config mentions an npm pin of \texttt{1.14.7}
        that was not runtime-enforced.
  \item \texttt{neuro\_pipeline} (v2.1.0) is an internal research workflow; it
        currently has no standalone peer-reviewed citation in the available
        materials and is therefore listed as software only.
  \item pydeface is cited as software (GitHub/Zenodo badge available upstream);
        no separate journal article DOI was verified here for the exact build
        \texttt{2.0.0+computecanada}.
\end{itemize}

\newpage

%==============================================================================
\section{Software inventory (versions used in this project)}
%==============================================================================

The table maps tools used in the curated workflow to the reference numbers in
Section~\ref{sec:refs}. Versions are those recorded in BIDS/derivatives metadata,
validation reports, or project documentation.

\begin{center}
\small
\begin{longtable}{@{}llp{0.42\textwidth}@{}}
\toprule
\textbf{Tool / resource} & \textbf{Version / identifier used} & \textbf{Role in this dataset} \\
\midrule
\endfirsthead
\toprule
\textbf{Tool / resource} & \textbf{Version / identifier used} & \textbf{Role in this dataset} \\
\midrule
\endhead
\bottomrule
\endfoot
BIDS specification & 1.9.0 & Dataset organisation target \\
\texttt{neuro\_pipeline} & 2.1.0 & Inventory, mapping, streaming de-id, conversion orchestration \\
dcm2niix & v1.0.20230411 & DICOM$\rightarrow$NIfTI + JSON sidecars \\
bids-validator & 1.15.0 (Scientific Data validation run) & BIDS validation \\
MRIQC & 24.1.0.dev0+gd5b13cb5.d20240826 (derivatives metadata); container target 24.0.2 in array scripts & Automated IQMs / QC reports \\
Pizarro model & \texttt{model.FINAL.onnx} (\texttt{Pizarro2023-FINAL}) & Complementary T1w artifact screening \\
ONNX Runtime & 1.23.2 & Pizarro inference engine \\
Python (Pizarro env) & 3.11.5 & Pizarro screening runtime \\
Monte Carlo dropout & 10 runs; seed 1010 & Pizarro uncertainty/probability estimation \\
pydicom & 3.0.2+computecanada & DICOM de-identification I/O \\
\texttt{mri\_anonymization} & 2.0.0 & De-identification engine \\
pydeface & 2.0.0+computecanada & Anatomical defacing \\
FSL & 6.0.7.20 & Dependency for pydeface registration \\
DICOM PS3.15 & confidentiality profile guidance & De-identification policy inspiration (not claimed as full-profile certification) \\
\end{longtable}
\end{center}

%==============================================================================
\section{Suggested in-text citation map}
%==============================================================================

\begin{itemize}[leftmargin=*]
  \item BIDS organisation: refs.~\ref{ref:bids}, optionally~\ref{ref:bidsapps}
  \item DICOM to NIfTI conversion (dcm2niix): ref.~\ref{ref:dcm2niix}
  \item MRIQC: ref.~\ref{ref:mriqc}
  \item Pizarro screening: ref.~\ref{ref:pizarro}
  \item FSL (pydeface dependency): ref.~\ref{ref:fsl}
  \item DICOM confidentiality guidance: ref.~\ref{ref:dicomps315}
  \item Optional sharing context (OpenNeuro): ref.~\ref{ref:openneuro}
\end{itemize}

%==============================================================================
\section{References}
\label{sec:refs}
%==============================================================================

\begin{enumerate}[leftmargin=*,itemsep=0.65em,label={[\arabic*]}]

  \item \label{ref:bids}
  Gorgolewski, K.~J. \textit{et al.}
  The brain imaging data structure, a format for organizing and describing
  outputs of neuroimaging experiments.
  \textit{Sci.\ Data} \textbf{3}, 160044 (2016).
  \url{https://doi.org/10.1038/sdata.2016.44}

  \item \label{ref:bidsapps}
  Gorgolewski, K.~J. \textit{et al.}
  BIDS apps: Improving ease of use, accessibility, and reproducibility of
  neuroimaging data analysis methods.
  \textit{PLoS Comput.\ Biol.} \textbf{13}, e1005209 (2017).
  \url{https://doi.org/10.1371/journal.pcbi.1005209}

  \item \label{ref:dcm2niix}
  Li, X., Morgan, P.~S., Ashburner, J., Smith, J.\ \& Rorden, C.
  The first step for neuroimaging data analysis: DICOM to NIfTI conversion.
  \textit{J.\ Neurosci.\ Methods} \textbf{264}, 47--56 (2016).
  \url{https://doi.org/10.1016/j.jneumeth.2016.03.001}

  \item \label{ref:mriqc}
  Esteban, O. \textit{et al.}
  MRIQC: Advancing the automatic prediction of image quality in MRI from unseen
  sites.
  \textit{PLoS ONE} \textbf{12}, e0184661 (2017).
  \url{https://doi.org/10.1371/journal.pone.0184661}

  \item \label{ref:pizarro}
  Pizarro, R. \textit{et al.}
  Deep learning, data ramping, and uncertainty estimation for detecting
  artifacts in large, imbalanced databases of MRI images.
  \textit{Med.\ Image Anal.} \textbf{90}, 102942 (2023).
  \url{https://doi.org/10.1016/j.media.2023.102942}\\
  Software implementation used here:
  \url{https://github.com/AS-Lab/Pizarro-et-al-2023-DL-detects-MRI-artifacts}
  (production model file \texttt{model.FINAL.onnx}).

  \item \label{ref:fsl}
  Smith, S.~M. \textit{et al.}
  Advances in functional and structural MR image analysis and implementation as
  FSL.
  \textit{NeuroImage} \textbf{23}, S208--S219 (2004).
  \url{https://doi.org/10.1016/j.neuroimage.2004.07.051}

  \item \label{ref:openneuro}
  Markiewicz, C.~J. \textit{et al.}
  The OpenNeuro resource for sharing of neuroscience data.
  \textit{eLife} \textbf{10}, e71774 (2021).
  \url{https://doi.org/10.7554/eLife.71774}

  \item \label{ref:numpy}
  Harris, C.~R. \textit{et al.}
  Array programming with NumPy.
  \textit{Nature} \textbf{585}, 357--362 (2020).
  \url{https://doi.org/10.1038/s41586-020-2649-2}

  \item \label{ref:matplotlib}
  Hunter, J.~D.
  Matplotlib: A 2D graphics environment.
  \textit{Comput.\ Sci.\ Eng.} \textbf{9}, 90--95 (2007).
  \url{https://doi.org/10.1109/MCSE.2007.55}

  \item \label{ref:dicomps315}
  National Electrical Manufacturers Association (NEMA).
  \textit{Digital Imaging and Communications in Medicine (DICOM) Standard,
  Part 15: Security and System Management Profiles}
  (Attribute Confidentiality Profiles).
  Rosslyn, VA: NEMA.
  Available at:
  \url{https://dicom.nema.org/medical/dicom/current/output/html/part15.html}
  (accessed for policy guidance; this dataset documents a custom workflow
  inspired by PS3.15 recommendations rather than certified full-profile
  compliance).

\end{enumerate}

%==============================================================================
\section{Software and standards (versioned resources)}
%==============================================================================

The following resources were used and should be named with version in Methods.
They are listed here as software/standards rather than journal articles when no
verified paper citation was required for this draft, or when the canonical
reference is the software distribution itself.

\begin{enumerate}[leftmargin=*,itemsep=0.55em,label={[S\arabic*]}]

  \item \textbf{bids-validator}
  (v1.15.0 used in Scientific Data validation reports).
  \url{https://github.com/bids-standard/bids-validator}

  \item \textbf{ONNX Runtime}
  (v1.23.2 used for Pizarro inference).
  \url{https://onnxruntime.ai/}

  \item \textbf{pydeface}
  (v2.0.0+computecanada used for anatomical defacing).
  \url{https://github.com/poldracklab/pydeface}

  \item \textbf{pydicom}
  (v3.0.2+computecanada observed in de-identification runtime).
  \url{https://github.com/pydicom/pydicom}

  \item \textbf{nibabel}
  (used for geometry checks in defacing QC workflows).
  \url{https://nipy.org/nibabel/}

  \item \textbf{Python}
  (v3.11.5 in the Pizarro screening environment).
  \url{https://www.python.org/}

  \item \textbf{neuro\_pipeline}
  (v2.1.0; internal conversion/de-identification/validation orchestration used
  to generate this BIDS dataset).
  Project software; no standalone peer-reviewed citation was available in the
  materials used for this draft.

\end{enumerate}

%==============================================================================
\section{Pizarro-specific citation note for Methods}
%==============================================================================

Suggested Methods wording (already aligned with the Scientific Data Pizarro
package):

\begin{quote}
\textit{An automated deep-learning classifier
(Pizarro et~al., 2023; model file \texttt{model.FINAL.onnx}) was applied to
T1-weighted images as a complementary quality-screening tool using ONNX Runtime.
Monte Carlo dropout inference (10 runs; seed 1010) provided continuous
model-estimated artifact probability, confidence, and uncertainty scores used
to prioritize targeted visual inspection. Model outputs were not used for
exclusion decisions. Quantitative MRI quality assessment remains provided by
MRIQC (Esteban et~al., 2017).}
\end{quote}

Corresponding reference entries: [\ref{ref:pizarro}] and [\ref{ref:mriqc}].

\vspace{1.5em}
\noindent\textit{End of References draft.}

\end{document}
